\begin{figure}[ht]
\centering
\begin{tikzpicture}[x=1.0cm, y=1.0cm, font=\small, >=Stealth]

% A plane drawn as a parallelogram (the subspace W = col(A))
\fill[engblue!10] (-0.5, 0) -- (5.0, 0) -- (6.5, 1.6) -- (1.0, 1.6) -- cycle;
\draw[engblue, thick] (-0.5, 0) -- (5.0, 0) -- (6.5, 1.6) -- (1.0, 1.6) -- cycle;
\node[engblue, anchor=west, font=\scriptsize] at (5.2, 0.3) {$W = \col(\mat{A})$};

% Origin
\fill[engblack] (1.5, 0.55) circle (1.4pt);
\node[anchor=north east, font=\scriptsize] at (1.5, 0.55) {$\vect{0}$};

% Point b above the plane
\fill[engred] (3.7, 3.0) circle (1.6pt);
\node[engred, anchor=west, font=\scriptsize] at (3.85, 3.05) {$\vect{b}$};

% Foot of the perpendicular (projection)
\fill[engamber] (3.7, 0.85) circle (1.6pt);
\node[engamber, anchor=north, font=\scriptsize] at (3.7, 0.72) {$\proj_{W}\vect{b}$};

% Residual (perpendicular) segment
\draw[engred, thick, ->] (3.7, 0.85) -- (3.7, 2.92);
\node[engred, anchor=west, font=\scriptsize] at (3.78, 1.9) {$\vect{b} - \proj_{W}\vect{b}$};

% Right-angle mark at the foot
\draw[engamber] (3.7, 0.85) ++(-0.28, 0) -- ++(0, 0.28) -- ++(0.28, 0);

% Vector from origin to b and to the projection
\draw[enggray, dashed] (1.5, 0.55) -- (3.7, 3.0);
\draw[enggray, dashed] (1.5, 0.55) -- (3.7, 0.85);

\node[anchor=north, font=\scriptsize, text width=11cm, align=center]
  at (2.8, -0.6) {
  The projection is the closest point of the subspace to $\vect{b}$, and
  the residual meets the subspace at a right angle. The length of the
  residual is the point-to-subspace distance the least-squares solve
  reports.
};

\end{tikzpicture}
\caption[Distance from a point to a subspace by projection]{The
orthogonal projection $\proj_{W}\vect{b}$ is the foot of the
perpendicular dropped from $\vect{b}$ to the subspace $W$, and the
residual $\vect{b} - \proj_{W}\vect{b}$ is orthogonal to every vector
in $W$. Among all points of $W$, the projection is the unique nearest
one to $\vect{b}$, so its distance $\norm{\vect{b} - \proj_{W}\vect{b}}_{2}$
is the least-squares residual. Reading the least-squares solution as a
distance-minimising projection is the geometric content behind the
normal equations.}
\label{fig:vol02:ch09:point-plane-distance}
\end{figure}
